\section{Datasets}

\paragraph{HBN development data.}
The Healthy Brain Network is a pediatric resource with EEG and other
measurements \citep{alexander2017hbn}. We used resting-state recordings.
Heads were fitted on a fixed HBN training
split, and checkpoints were selected on a separate validation split. The
previously inspected HBN R5 test split was excluded from both steps.
The source manifest identifies 800 training participants and 100 validation
participants, with one recording per participant. Mean age was
$10.17 \pm 3.38$ years in training (range 5.06--21.67) and
$10.17 \pm 3.34$ years in validation (range 5.06--21.18), where the spread is
the sample standard deviation. Validation used release R8; training used
R1--R4, R6--R7, and R9--R10. The manifest's SHA-256 matches the source identity
retained in the primary analysis. The 200-, 400-, and 800-subject subsets
described below belong to the secondary training-size analysis.

\paragraph{MIPDB external evaluation.}
MIPDB contains high-density EEG recorded during tasks and rest across development
\citep{langer2017mipdb}. We used a fixed acquisition snapshot from NEMAR
\texttt{nm000153}. Of 126 candidate records, 17 had no EEG recording. The
remaining 109 participants were assigned to a 10-person engineering pilot,
79 primary candidates within HBN's training-age range, and 20 participants
outside that range. Four primary candidates failed the prespecified signal
quality checks, leaving \PrimarySubjects{} participants for the main analysis.
The pilot was used only to check the recording pipeline, without age-prediction
metrics. The out-of-range group was excluded from the primary comparison.

\begin{table}[tbp]
  \centering
  \caption{MIPDB groups after allocation and signal quality checks.}
  \label{tab:cohorts}
  % Generated table; do not edit.
\begin{tabular}{lrl}
\toprule
Cohort & Subjects & Role \\
\midrule
MIPDB pilot & 10 & Protocol/QC pilot \\
MIPDB primary & 75 & Sealed confirmatory cohort \\
MIPDB extrapolation & 20 & Descriptive boundary cohort \\
\bottomrule
\end{tabular}

\end{table}

MIPDB predictions and targets were not used to select the primary heads,
checkpoints, or decision thresholds. The later training-size and layer-wise
analyses reused these participants, so they are exploratory follow-ups rather
than independent replications.

\paragraph{Separate transfer cohort.}
The SFARI\_EEG multi-paradigm dataset, OpenNeuro ds006780 version 1.0.0,
contains BioSemi-64 EEG recordings from children in typical-development, autism,
and sibling groups \citep{molholmsfari}. We selected \texttt{Restingstate}
run-01. Of 129 matching recordings, one lacked the required events file.
All 128 remaining recordings passed the documented signal checks. One
participant had no metadata row and another had no age, leaving
\DsSubjectCount{} participants. Their mean age was $10.60 \pm 1.57$ years
(range 7.9--14.5); 80 were recorded as male and 46 as female. The source group
labels identified 39 typically developing participants, 62 with autism, and
25 in the sibling group. These descriptions were recovered from the pinned
public metadata: hashes of both the ordered participant list and age vector
exactly match the retained v5 analysis. Family relationships within this
sample remain unknown.

\paragraph{Data use.}
This study reuses existing datasets and involved no new participant recruitment.
The HBN and MIPDB source publications describe their original ethics and
consent procedures \citep{alexander2017hbn,langer2017mipdb}. The ds006780
metadata report approval by the Albert Einstein College of Medicine
Institutional Review Board (2021-13433), participant assent, parental or guardian
consent, and a CC0 data license \citep{molholmsfari}. These statements concern
the original data collection. The analysis bundles contain aggregate results;
access to source recordings remains subject
to the respective providers' terms.
